\problemname{DFA Minimzation}

\begin{quote}
	This is an example of a \enquote{direct} task, where the main challenge simply is implementing a given algorithm.
\end{quote}

Camila is currently studying for her automata theory exam and is worried about DFA minimization.
She has solved several exercises on paper, but wants to double check that she didn't make a mistake.
For this, just knowing the number of states in the minimized automaton is already very helpful for her.
Can you help her?

\begin{Input}
The input begins with a line consisting of three space-separated integers $n$, the number of states in the automaton, numbered from $1$ to $n$, $m$, the size of the alphabet, and $i$, the initial state.
On the next line, an integer $k$ gives the number of accepting states, followed by $k$ space separated integers, denoting the accepting states.
Finally, $n$ lines follow, describing the transitions of the automaton.
The $i$-th line comprises $m$ space separated integers.
The $j$-th integer describes the successor of $i$ under the $j$-th letter.
\end{Input}

\begin{Output}
Output an integer $n$ where $n$ is the number of states in the minimal DFA.
\end{Output}

\begin{Constraints}
\begin{itemize}
	\item $1 \leq n \leq 1000$
	\item $1 \leq m \leq 24$
	\item $0 \leq k \leq n$
	\item Each reference to a state (e.g.\ in the transition matrix) is well-formed, i.e.\ between $1$ and $n$.
\end{itemize}
\end{Constraints}
